\documentclass[12pt, a4paper]{article}
\usepackage[UTF8]{ctex}
\usepackage{amsmath, amssymb, amsthm}
\usepackage{geometry}
\usepackage{bm}

\geometry{left=2.5cm, right=2.5cm, top=2.5cm, bottom=2.5cm}

\newcommand{\RR}{\mathbb{R}}
\newcommand{\e}{\mathrm{e}}
\newcommand{\dif}{\mathrm{d}}

\begin{document}
\section*{15.2}
\textbf{4.}讨论下列含参变量反常积分的一致收敛性：

(2)\(\int_{-\infty}^{+\infty} \e^{-(x - \alpha)^2}\dif x\)，在(i)\(a < \alpha < b\)；(i)\(-\infty < \alpha < +\infty\).

(i)考虑尾部积分
\[\int_{A}^{+\infty} \e^{-(x - \alpha)^2}\dif x + \int_{-\infty}^{-A} \e^{-(x - \alpha)^2}\dif x.\]
而
\[\int_{A}^{+\infty} \e^{-(x - \alpha)^2}\dif x = \int_{A - \alpha}^{+\infty} \e^{-u^2}\dif u \leq \int_{A - b}^{+\infty} \e^{-u^2}\dif u \to 0,\]
同理
\[\int_{-\infty}^{-A} \e^{-(x - \alpha)^2}\dif x \leq \int_{-\infty}^{-A - a} \e^{-u^2}\dif u \to 0,\]
故原积分一致收敛.

(ii)令\(\alpha = A\)，于是
\[\int_{A}^{+\infty} \e^{-(x - \alpha)^2}\dif x = \int_{A - \alpha}^{+\infty} \e^{-u^2}\dif u = \int_{0}^{+\infty} \e^{-u^2}\dif u = \frac{\sqrt{\pi}}{2},\]
故原积分不一致收敛.

(4)\(\int_{0}^{+\infty} \e^{-\alpha x}\sin x\dif x\)，在(i)\(\alpha \geq \alpha_0 > 0\)；(ii)\(\alpha > 0\).

(i)显然\(\int_0^A \sin x\dif x \forall A\)一致有界，\(\e^{-\alpha x}\)单调且\(x \to +\infty, \e^{-\alpha x} \leq \e^{-\alpha_0 x} \to 0\).故原积分一致收敛.

(ii)考虑尾部积分
\[\int_{A}^{+\infty} \e^{-\alpha x}\sin x\dif x,\]
取\(\alpha = \dfrac{1}{A}\)，当\(\alpha \to 0+\)时有
\[\lim_{\alpha \to 0+}\int_{A}^{+\infty} \e^{-\alpha x}\sin x\dif x = \int_{A}^{+\infty} \sin x\dif x.\]
故原积分不一致收敛.

\textbf{6.}确定函数\(F(y) = \int_{0}^{\pi} \dfrac{\sin x}{x^y(\pi - x)^{2 - y}}\dif x\)的连续范围.

解：
\begin{align*}
F(y) &= \int_{0}^{\pi} \frac{\sin x}{x^y(\pi - x)^{2 - y}}\dif x \\
&= \int_0^\varepsilon \frac{\sin x}{x^y(\pi - x)^{2 - y}}\dif x + \int_{\varepsilon}^{\pi - \varepsilon} \frac{\sin x}{x^y(\pi - x)^{2 - y}}\dif x + \int_{\pi - \varepsilon}^{\pi} \frac{\sin x}{x^y(\pi - x)^{2 - y}}\dif x \\
\end{align*}
第一个积分在\(1 - y > -1\)，即\(y < 2\)时收敛，第二个积分收敛，第三个积分在\(y - 1 > -1\)，即\(y > 0\)时收敛，故\(F(y)\)在\(0 < y < 2\)连续.

\textbf{8.}证明：函数\(I(t) = \int_{0}^{+\infty} \dfrac{\cos x}{1 + (x + t)^2}\dif x\)在\((-\infty, +\infty)\)上可微.

解：\(f(x, t) = \dfrac{\cos x}{1 + (x + t)^2}\)，
\[\int_{0}^{+\infty} |f(x, t)|\dif x \leq \int_{0}^{+\infty} \frac{1}{1 + (x + t)^2}\dif x = \int_{t}^{+\infty} \frac{1}{1 + u^2}\dif u \leq \frac{\pi}{2}.\]
\(\forall [a, b] \subset (-\infty, +\infty), \forall t \in [a, b]\)，设\(C = \max\{|a|, |b|\}\)，
\[|f_t(x, t)| = \left|-\frac{2(x + t)}{(1 + (x + t)^2)^2}\cos x\right| \leq \frac{2|x + t|}{(1 + (x + t)^2)^2} \leq \frac{2(x + C)}{\left(1 + \dfrac{x^2}{4}\right)^2}.\]
故\(\exists M > 0, X_0 > 0\)，使当\(x \geq X_0\)时，\(\forall t \in [a, b]\)有
\[|f_t(x, t)| \leq \frac{M}{x^3},\]
且\(\int_{X_0}^{+\infty} \frac{M}{x^3}\dif x\)收敛，故\(\int_{0}^{+\infty} f_(x, t)\dif x\)在\(\RR\)内闭一致收敛，故\(I(t)\)在\(\RR\)上可微.

\textbf{9.}利用\(\dfrac{\e^{-ax} - \e^{-bx}}{x} = \int_a^b \e^{-xy}\dif y\).计算\(\int_{0}^{+\infty} \dfrac{\e^{-ax} - \e^{-bx}}{x}\dif x \quad (b > a > 0)\).

解：
\begin{align*}
\int_{0}^{+\infty} \frac{\e^{-ax} - \e^{-bx}}{x}\dif x &= \int_{0}^{+\infty} \dif x\int_a^b \e^{-xy}\dif y \\
&= \int_a^b \dif y\int_{0}^{+\infty} \e^{-xy}\dif x \\
&= \ln\frac{b}{a}.
\end{align*}

\textbf{10.}利用\(\dfrac{\sin bx - \sin ax}{x} = \int_a^b \cos xy\dif y\)，计算\(\int_{0}^{+\infty} \e^{-px}\dfrac{\sin bx - \sin ax}{x}\dif x\).

解：
\begin{align*}
\int_{0}^{+\infty} \e^{-px}\frac{\sin bx - \sin ax}{x}\dif x &= \int_{0}^{+\infty} \e^{-px}\dif x\int_a^b \cos xy\dif y \\
&= \int_a^b \dif y\int_{0}^{+\infty} \e^{-px} \cos xy\dif x \\
&= \int_a^b \frac{p}{p^2 + y^2}\dif y \\
&= \arctan\frac{b}{p} - \arctan\frac{a}{p}.
\end{align*}

\textbf{13.}设\(f(x)\)在\([0, +\infty)\)上连续，且\(\lim\limits_{x \to +\infty} f(x) = 0\)，证明
\[\int_{0}^{+\infty} \frac{f(ax) - f(bx)}{x}\dif x = f(0)\ln\frac{b}{a} \quad (a, b > 0).\]

解：
\begin{align*}
\int_{0}^{+\infty} \frac{f(ax) - f(bx)}{x}\dif x &= \int_{0}^{+\infty} \dif x\int_b^a f'(xy) \dif y \\
&= \int_b^a \dif y\int_{0}^{+\infty} f'(xy) \dif x \\
&= \int_b^a \frac{1}{y} \dif y f(0) \\
&= f(0)\ln\frac{b}{a}.
\end{align*}

\textbf{14.}(1)利用\(\int_{0}^{+\infty} \e^{-y^2}\dif y = \frac{\sqrt{\pi}}{2}\)推出\(L(c) = \int_{0}^{+\infty} \e^{-y^2 - \frac{c^2}{y^2}}\dif y = \frac{\sqrt{\pi}}{2}\e^{-2c} \quad (c > 0)\)；

(2)利用积分号下求导的方法引出\(\dfrac{\dif L}{\dif c} = -2L\)，以此推出与(1)同样的结果，并计算\(\int_{0}^{+\infty} \e^{-ay^2 - \frac{b}{y^2}}\dif y \quad (a > 0, b > 0)\).

解：(1)令\(t = \dfrac{c}{y}\)
\begin{align*}
L(c) &= \int_{0}^{+\infty} \e^{-y^2 - \frac{c^2}{y^2}}\dif y \\
&= \int_{0}^{+\infty} \frac{c}{t^2}\e^{-\frac{c^2}{t^2} - t^2}\dif t.
\end{align*}
\[2L(c) = \int_{0}^{+\infty} \left(1 + \frac{c}{y^2}\right)\e^{-y^2 - \frac{c^2}{y^2}}\dif y.\]
令\(u = y - \dfrac{c}{y}, \dif u = \left(1 + \dfrac{c}{y^2}\right)\dif y\)，有
\[L(c) = \frac{1}{2}\int_{-\infty}^{+\infty} \e^{-u^2 - 2c}\dif u = \frac{\sqrt{\pi}}{2}\e^{-2c}.\]

(2)
\begin{align*}
\frac{\dif L}{\dif c} &= \int_{0}^{+\infty} \frac{\partial}{\partial c}\left(\e^{-y^2 - \frac{c^2}{y^2}}\right)\dif y \\
&= \int_{0}^{+\infty} \left(-\frac{2c}{y^2}\right)\e^{-y^2 - \frac{c^2}{y^2}}\dif y \\
&= \int_{+\infty}^{0} \left(-\frac{2c}{y^2}\right)\left(-\frac{c}{t^2}\right)\e^{-y^2 - \frac{c^2}{y^2}}\dif y \\
&= -2L
\end{align*}
于是\(L(c) = A\e^{-2c}\)，由\(L(0) = \dfrac{\sqrt{\pi}}{2}\)，得\(L(c) = \dfrac{\sqrt{\pi}}{2}\e^{-2c}\).

令\(t = \sqrt{a}y\)，则
\[\int_{0}^{+\infty} \e^{-ay^2 - \frac{b}{y^2}}\dif y = \frac{1}{\sqrt{a}}\int_{0}^{+\infty} \e^{-t^2 - \frac{ab}{t^2}}\dif t = \frac{\sqrt{\pi}}{2\sqrt{a}}\e^{-2\sqrt{ab}}.\]

\textbf{15.}利用\(\int_{0}^{+\infty} \e^{-t(\alpha^2 + x^2)}\dif t = \dfrac{1}{\alpha^2 + x^2}\)，计算\(J = \int_{0}^{+\infty} \dfrac{\cos\beta x}{\alpha^2 + x^2}\dif x \quad (\alpha > 0)\).

解：
\begin{align*}
J &= \int_{0}^{+\infty} \frac{\cos\beta x}{\alpha^2 + x^2}\dif x \\
&= \int_{0}^{+\infty} \cos\beta x\dif x\int_{0}^{+\infty} \e^{-t(\alpha^2 + x^2)}\dif t \\
&= \int_{0}^{+\infty} \dif t\int_{0}^{+\infty} \e^{-t(\alpha^2 + x^2)}\cos\beta x\dif x \\
&= \int_{0}^{+\infty} \e^{-t\alpha^2}\dif t \int_{0}^{+\infty} \cos\beta x\e^{-tx^2}\dif x \\
&= \int_{0}^{+\infty} \frac{1}{2}\sqrt{\frac{\pi}{t}}\e^{-\frac{\beta^2}{4t}}\e^{-t\alpha^2}\dif t \\
&= \frac{\sqrt{\pi}}{2} \int_{0}^{+\infty} t^{-\frac{1}{2}}\e^{-\alpha^2t - \frac{\frac{\beta^2}{4}}{t}}\dif t \\
&= \frac{\pi}{2\alpha}\e^{-\alpha|\beta|}.
\end{align*}

\section*{15.3}
\textbf{1.}计算下列积分：

(2)\(\int_0^\pi \dfrac{\dif x}{\sqrt{3 - \cos x}}\).

\[I = \frac{1}{\sqrt{2}}\int_0^\pi \frac{\dif x}{\sqrt{1 + \sin^2\left(\dfrac{x}{2}\right)}} = \sqrt{2}\int_{0}^{\frac{\pi}{2}} \frac{\dif t}{\sqrt{1 + \sin^2t}} = \sqrt{2}K(i) = \frac{\Gamma^2\left(\frac{1}{4}\right)}{4\sqrt{\pi}}.\]

(4)\(\int_{0}^{+\infty} \dfrac{x^{m - 1}}{1 + x^n}\dif x \quad (n > m > 0)\).

令\(x = t^{\frac{1}{n}}\)，
\[I = \int_{0}^{+\infty} \frac{t^{\frac{m - 1}{n}}}{1 + t} \cdot \frac{1}{n}t^{\frac{1}{n - 1}}\dif t = \frac{1}{n}\int_{0}^{+\infty} \frac{t^{\frac{m}{n} - 1}}{1 + t}\dif t = \frac{\pi}{n\sin\dfrac{m\pi}{n}}.\]

(6)\(\int_{0}^{\frac{\pi}{2}} \sin^7x\cos^{\frac{1}{2}}x\dif x\).

\[I = \frac{1}{2}B\left(4, \frac{3}{4}\right) = \frac{\Gamma(4)\Gamma\left(\dfrac{3}{4}\right)}{2\Gamma\left(\dfrac{19}{4}\right)} = \frac{256}{1155}.\]

(8)\(\int_0^1 x^{p - 1}(1 - x^n)^{q - 1}\dif x \quad (p, q, n > 0)\).

令\(x = t^{\frac{1}{n}}\)，
\[I = \frac{1}{n}\int_0^1 t^{\frac{p}{n} - 1}(1 - t)^{q - 1}\dif t = \frac{B\left(\dfrac{p}{n}, q\right)}{n}.\]

\textbf{2.}证明\(\int_{0}^{+\infty} \e^{-x^n}\dif x = \dfrac{1}{n}\Gamma\left(\dfrac{1}{n}\right)\)（\(n\)为正整数），并推出\(\lim\limits_{n \to \infty} \int_{0}^{+\infty} \e^{-x^n}\dif x = 1\).

解：令\(x = t^{\frac{1}{n}}\)，
\[\int_{0}^{+\infty} \e^{-x^n}\dif x = \frac{1}{n}\int_{0}^{+\infty} t^{\frac{1}{n} - 1}\e^{-t}\dif t = \frac{1}{n}\Gamma\left(\frac{1}{n}\right).\]
\[\lim_{n \to \infty} \int_{0}^{+\infty} \e^{-x^n}\dif x = \lim_{n \to \infty} \frac{1}{n}\Gamma\left(\frac{1}{n}\right) = \lim_{n \to \infty} \Gamma(1 + \frac{1}{n}) = 1.\]

\textbf{4.}证明\(\lim\limits_{s \to +\infty} \Gamma(s) = +\infty\).

解：
\[\Gamma(s) = \int_{0}^{+\infty} x^{s - 1}\e^{-x} \dif x \geq \int_{a}^{+\infty} x^{s - 1}\e^{-x}\dif x \geq a^{s - 1}\int_{a}^{+\infty}\e^{-x}\dif x = a^{s - 1}\e^{-a} \to +\infty \quad (s \to +\infty), \forall a > 1.\]

\textbf{5.}计算\(\int_0^1 \ln\Gamma(x)\dif x\).

解：
\begin{align*}
\Gamma(x)\Gamma(1 - x) &= \frac{\pi}{\sin\pi x} \\
\ln\Gamma(x) + \ln\Gamma(1 - x) &= \ln\pi - \ln\sin\pi x \\
\int_0^1 \ln\Gamma(x)\dif x + \int_0^1 \ln\Gamma(1 - x)\dif x &= \int_0^1 (\ln\pi - \ln\sin\pi x)\dif x \\
2\int_0^1 \ln\Gamma(x)\dif x &= \ln\pi - \int_0^1 \ln\sin\pi x\dif x \\
\int_0^1 \ln\Gamma(x)\dif x = \ln\sqrt{2\pi}
\end{align*}

\textbf{6.}设\(\Omega = \{(x, y, z)|x^2 + y^2 + z^2 \leq 1\}\).确定正数\(p\)，使得反常重积分
\[I = \iiint\limits_\Omega \frac{\dif x\dif y\dif z}{(1 - x^2 - y^2 - z^2)^p}\]
收敛.并在收敛时计算\(I\)的值.

解：设\(x^2 + y^2 + z^2 = r^2\)，则
\[I = 4\pi\int_0^1 \frac{r^2}{(1 - r^2)^p}\dif r,\]
当且仅当\(p < 1\)时收敛.令\(t = r^2\)，
\[I = 2\pi B\left(\frac{3}{2}, 1 - p\right).\]

\textbf{7.}设\(\Omega = \{(x, y, z)|x \geq 0, y \geq 0, z \geq 0\}\).确定正数\(\alpha, \beta, \gamma\)，使得反常重积分
\[I = \iiint\limits_\Omega \frac{\dif x\dif y\dif z}{1 + x^\alpha + y^\beta + z^\gamma}\]
收敛，并在收敛时计算\(I\)的值.

解：令
\[\begin{cases}
x = u^{\frac{2}{\alpha}}, \\
y = v^{\frac{2}{\beta}}, \\
z = w^{\frac{2}{\gamma}},
\end{cases} \quad \begin{cases}
u = r\sin\varphi\cos\theta, \\
v = r\sin\varphi\sin\theta, \\
w = r\cos\varphi,
\end{cases}\]
得到
\begin{align*}
I &= \frac{8}{\alpha\beta\gamma}\int_{0}^{\frac{\pi}{2}} \sin^{\frac{2}{\beta} - 1}\theta\cos^{\frac{2}{\alpha} - 1}\theta\dif\theta \int_{0}^{\frac{\pi}{2}} \sin^{\frac{2}{\alpha} + \frac{2}{\beta} - 1}\varphi\cos^{\frac{2}{\gamma} - 1}\varphi\dif\varphi \int_{0}^{+\infty} \frac{r^{\frac{2}{\alpha} + \frac{2}{\beta} + \frac{2}{\gamma} - 1}}{1 + r^2}\dif r \\
&= \frac{1}{\alpha\beta\gamma}\frac{\Gamma\left(\dfrac{1}{\alpha}\right)\Gamma\left(\dfrac{1}{\beta}\right)}{\Gamma\left(\dfrac{1}{\alpha} + \dfrac{1}{\beta}\right)}\frac{\Gamma\left(\dfrac{1}{\alpha} + \dfrac{1}{\beta}\right)\Gamma\left(\dfrac{1}{\gamma}\right)}{\Gamma(p)}\frac{\pi}{\sin(\pi p)} \\
&= \frac{1}{\alpha\beta\gamma}\Gamma\left(\frac{1}{\alpha}\right)\Gamma\left(\frac{1}{\beta}\right)\Gamma\left(\frac{1}{\gamma}\right)\Gamma\left(1 - \frac{1}{alpha} - \frac{1}{\beta} - \frac{1}{\gamma}\right).
\end{align*}
当且仅当\(\dfrac{1}{\alpha} + \dfrac{1}{\beta} + \dfrac{1}{\gamma} < 1\)时收敛.

\textbf{8.}计算
\[I = \iint\limits_D x^{m - 1}y^{n - 1}(1 - x - y)^{p - 1}\dif x\dif y,\]
其中\(D\)是由三条直线\(x = 0, y = 0\)及\(x + y = 1\)所围成的闭区域，\(m, n, p\)均为大于0的正数.

解：
\[I = (p - 1)\iiint\limits_\Omega x^{m - 1}y^{n - 1}z^{p - 2}\dif x\dif y\dif z,\]
其中\(\Omega\)为平面\(x = 0, y = 0, z = 0, x + y + z = 1\)所围成的区域.令
\[\begin{cases}
x = u^2, \\
y = v^2, \\
z = w^2,
\end{cases} \quad \begin{cases}
u = r\sin\varphi\cos\theta, \\
v = r\sin\varphi\sin\theta, \\
w = r\cos\varphi,
\end{cases}\]
得到
\[I = 8(p - 1)\int_{0}^{\frac{\pi}{2}} \sin^{2n - 1}\theta\cos^{2m - 1}\theta\dif\theta \int_{0}^{\frac{\pi}{2}}\sin^{2m + 2n - 1}\varphi\cos^{2m - 1}\varphi\dif\varphi \int_0^1 r^{2m + 2n + 2p - 3}\dif r = \frac{\Gamma(m)\Gamma(n)\Gamma(p)}{\Gamma(m + n + p)}.\]

\textbf{10.}设\(0 < \alpha < 2, 0 < k < 1\)，证明：
\[\int_0^\pi \left(\frac{\sin\varphi}{1 + \cos\varphi}\right)^{\alpha - 1}\frac{\dif\varphi}{1 + k\cos\varphi} = \frac{1}{1 + k}\left(\sqrt{\frac{1 + k}{1 - k}}\right)^\alpha\frac{\pi}{\sin\dfrac{\alpha}{2}\pi}.\]

解：令\(t = \tan\dfrac{\varphi}{2}\)，则
\[\int_0^\pi \left(\frac{\sin\varphi}{1 + \cos\varphi}\right)^{\alpha - 1}\frac{\dif\varphi}{1 + k\cos\varphi} = 2\int_{0}^{+\infty} \frac{t^{\alpha - 1}\dif t}{(1 + k) + (1 - k)t^2},\]
令\(\sqrt{\dfrac{1 - k}{1 + k}}t = \tan\theta\)，则
\begin{align*}
2\int_{0}^{+\infty} \frac{t^{\alpha - 1}\dif t}{(1 + k) + (1 - k)t^2} &= \frac{2}{1 + k}\left(\sqrt{\frac{1 - k}{1 + k}}\right)^\alpha \int_{0}^{\frac{\pi}{2}} \sin^{\alpha - 1}\theta\cos^{1 - \alpha}\theta\dif\theta \\
&= \frac{1}{1 + k}\left(\sqrt{\frac{1 - k}{1 + k}}\right)^\alpha B\left(\frac{\alpha}{2}, 1 - \frac{\alpha}{2}\right) \\
&= \frac{1}{1 + k}\left(\sqrt{\frac{1 - k}{1 + k}}\right)^\alpha \Gamma\left(\frac{\alpha}{2}\right)\Gamma\left(1 - \frac{\alpha}{2}\right) \\
&= \frac{1}{1 + k}\left(\sqrt{\frac{1 + k}{1 - k}}\right)^\alpha\frac{\pi}{\sin\dfrac{\alpha}{2}\pi}.
\end{align*}

\textbf{11.}设\(0 \leq h < 1\)，正整数\(n \geq 3\)，证明：
\[\int_0^h (1 - t^2)^{\frac{n - 3}{2}}\dif t \geq \frac{\sqrt{\pi}}{2}\frac{\Gamma\left(\dfrac{n - 1}{2}\right)}{\Gamma\left(\dfrac{n}{2}\right)}h.\]

解：令\(t = hu\)，得
\[\int_0^h (1 - t^2)^{\frac{n - 3}{2}}\dif t = h\int_0^1 (1 - h^2u^2)^{\frac{n - 3}{2}}\dif u \geq h\int_0^1 (1 - u^2)^{\frac{n - 3}{2}}\dif u.\]
令\(u = \sin\theta\)，则
\[h\int_0^1 (1 - u^2)^{\frac{n - 3}{2}}\dif u = h\int_{0}^{\frac{\pi}{2}} \cos^{n - 2}\theta\dif\theta = \frac{h}{2}B\left(\frac{1}{2}, \frac{n - 1}{2}\right) = \frac{h}{2}\frac{\Gamma\left(\dfrac{1}{2}\right)\Gamma\left(\dfrac{n - 1}{2}\right)}{\Gamma\left(\dfrac{n}{2}\right)} = \frac{\sqrt{\pi}}{2}\frac{\Gamma\left(\dfrac{n - 1}{2}\right)}{\Gamma\left(\dfrac{n}{2}\right)}h.\]

\end{document}